\documentclass[11pt]{article}

\usepackage{amsmath, amssymb, amsfonts}
\usepackage{geometry}
\geometry{margin=1in}

\title{Supplementary Information:\\
Bayesian Evidence Correction via the\\
Posterior Information Distortion}
\author{}
\date{}

\begin{document}

\maketitle

\begin{abstract}
\noindent We derive the correction to the log-evidence that follows from
replacing the naive Bayesian Laplace covariance $\mathbf H_{\mathrm{post}}^{-1}$
with the sandwich-corrected posterior covariance
$\Sigma_{\mathrm{DCC},\mathrm{post}} = \mathbf H_{\mathrm{post}}^{-1}\,
\mathbf{JT}_{\mathrm{post}}\,\mathbf H_{\mathrm{post}}^{-1}$ in the Laplace
approximation. The result is the direct posterior counterpart of the likelihood
evidence correction of \texttt{supplement\_evidence\_correction.tex}: the shift
in log-evidence equals $\tfrac{1}{2}\log\det(\mathbf C_{\mathrm{post}})$, where
$\mathbf C_{\mathrm{post}}$ is the Posterior Information Distortion (PID). At
the population level under correct likelihood specification, this shift vanishes
\emph{for any proper Gaussian prior covariance}, so the correction is in that
regime a pure likelihood-misspecification diagnostic. We further introduce an
\emph{evidence interval} framework that reports the naive,
$\mathbf{JT}_{\mathrm{post}}$-only, and DCC estimates of $\log Z$ jointly, their
pairwise differences being determined by $\log\det(\mathbf C_{\mathrm{post}})$.
\end{abstract}

\section*{1. Setup}

We adopt the notation of \texttt{supplement\_posterior\_main.tex}. Let $D = (y_1,\ldots,y_T)$
be the data, $\pi(\theta) = \mathcal{N}(\mu_{\mathrm{prior}}, \Sigma_{\mathrm{prior}})$ a
proper Gaussian prior with precision $\mathbf H_{\mathrm{prior}} = \Sigma_{\mathrm{prior}}^{-1}$,
and $\ell(\theta) = \log L(\theta)$ the log-likelihood. The posterior Hessian and total
posterior score fluctuation matrix decompose additively (Sections~2 and 3 of the main
supplement),
\begin{equation}
\mathbf H_{\mathrm{post}} \;=\; \mathbf H_{\mathrm{lik}} + \mathbf H_{\mathrm{prior}},
\qquad
\mathbf{JT}_{\mathrm{post}} \;:=\; \mathbf{JT}_{\mathrm{lik}} + \mathbf H_{\mathrm{prior}},
\label{eq:setup-decomp}
\end{equation}
with $\mathbf{JT}_{\mathrm{lik}} = \mathrm{Cov}(S_{\mathrm{lik}})$ the total likelihood
score covariance. The Posterior Information Distortion (PID) is
\begin{equation}
\mathbf C_{\mathrm{post}}
\;=\;
\mathbf H_{\mathrm{post}}^{-1/2}\,
\mathbf{JT}_{\mathrm{post}}\,
\mathbf H_{\mathrm{post}}^{-1/2},
\label{eq:Cpost-def}
\end{equation}
and the Distortion-Corrected posterior Covariance is
\begin{equation}
\Sigma_{\mathrm{DCC},\mathrm{post}}
\;=\;
\mathbf H_{\mathrm{post}}^{-1}\,
\mathbf{JT}_{\mathrm{post}}\,
\mathbf H_{\mathrm{post}}^{-1}
\;=\;
\mathbf H_{\mathrm{post}}^{-1/2}\,
\mathbf C_{\mathrm{post}}\,
\mathbf H_{\mathrm{post}}^{-1/2}.
\label{eq:SigmaDCCpost}
\end{equation}
At the population level under correctly specified likelihood, the information identity
$\mathbf{JT}_{\mathrm{lik}} = \mathbf H_{\mathrm{lik}}$ propagates additively to
$\mathbf{JT}_{\mathrm{post}} = \mathbf H_{\mathrm{post}}$ (Section~5.1 of the main
supplement), so $\mathbf C_{\mathrm{post}} = \mathbf I$ for any proper Gaussian
prior (Section~6.1 of the main supplement); see~\eqref{eq:Cpost-def}--\eqref{eq:SigmaDCCpost}.
The DCC then reduces to the standard Bayesian Laplace covariance
$\mathbf H_{\mathrm{post}}^{-1}$ (Section~9.4 of the main supplement).

\paragraph{Preconditions.}
All formulas below assume $\mathbf H_{\mathrm{post}} \succ 0$, which is automatic for
analytic $\mathbf H_{\mathrm{lik}} \succeq 0$ with any proper Gaussian prior, and otherwise
requires the dominance condition $\mathbf H_{\mathrm{prior}} - \mathbf H_{\mathrm{lik}}^- \succ 0$
of Section~2.1 of the main supplement (label \texttt{eq:dominance-condition} in that
supplement). The evidence correction $\Delta\log Z$ is not defined as a real number outside
this regime; in finite-difference Hessian implementations the dominance condition must be
checked.

\paragraph{Goal of this supplement.}
This supplement derives the Bayesian model-evidence correction obtained by replacing
the naive Bayesian Laplace covariance $\mathbf H_{\mathrm{post}}^{-1}$ with the
sandwich-corrected $\Sigma_{\mathrm{DCC},\mathrm{post}}$ in the Laplace approximation
to $\log Z = \log \int L(\theta)\,\pi(\theta)\,\mathrm{d}\theta$. The development
parallels \texttt{supplement\_evidence\_correction.tex} (the likelihood-side counterpart),
but the result has three distinctive features:
\begin{itemize}
\item the correction vanishes at the population level under correct likelihood
specification \emph{for any proper Gaussian prior covariance}, exhibiting a
prior-cancellation property absent from the likelihood-side correction in the
singular-Fisher regime;
\item the three natural Laplace estimators (naive, $\mathbf{JT}_{\mathrm{post}}$-only,
and DCC) form an interval on $\log Z$ whose width is exactly $|\log\det(\mathbf C_{\mathrm{post}})|$
and which we report as a calibration uncertainty;
\item the correction is an $O(1)$ intensive quantity, decisive only in the regime where
the $O(T)$ likelihood values cancel \emph{and} the prior is not strong enough to suppress
the likelihood-side residual.
\end{itemize}

\section*{2. Laplace Approximation to the Bayesian Evidence}

Let $Z := \int L(\theta)\,\pi(\theta)\,\mathrm{d}\theta$ be the Bayesian evidence. We expand
the integrand around the maximum-a-posteriori (MAP) estimate
$\hat\theta := \arg\max_\theta [\log L(\theta) + \log\pi(\theta)]$,
\begin{equation}
\log L(\theta) + \log\pi(\theta)
\;\approx\;
\log L(\hat\theta) + \log\pi(\hat\theta)
\;-\;
\tfrac{1}{2}\,(\theta - \hat\theta)^\top
\boldsymbol\Lambda\,
(\theta - \hat\theta),
\label{eq:quadratic-expansion}
\end{equation}
where $\boldsymbol\Lambda$ is the precision matrix of the local Gaussian approximation
to the posterior at $\hat\theta$. Gaussian integration yields
\begin{equation}
\log Z
\;\approx\;
\log L(\hat\theta) + \log\pi(\hat\theta)
\;+\;
\tfrac{p}{2}\log(2\pi)
\;-\;
\tfrac{1}{2}\log\det(\boldsymbol\Lambda),
\label{eq:laplace-precision}
\end{equation}
with $p = \dim(\theta)$. Equivalently, in terms of the posterior covariance
$\Sigma = \boldsymbol\Lambda^{-1}$:
\begin{equation}
\log Z
\;\approx\;
\log L(\hat\theta) + \log\pi(\hat\theta)
\;+\;
\tfrac{p}{2}\log(2\pi)
\;+\;
\tfrac{1}{2}\log\det(\Sigma).
\label{eq:laplace-covariance}
\end{equation}
Equations \eqref{eq:laplace-precision} and \eqref{eq:laplace-covariance} are exact
representations of the same approximation, written in precision and covariance form;
the choice of which to instantiate determines which curvature object is identified with
the posterior precision.

\paragraph{Remark on Jacobians.}
As in the main supplement, all derivatives are taken in the transformed coordinates in
which the prior is Gaussian, so the integration in $Z$ is over the transformed parameter
and no Jacobian corrections appear in \eqref{eq:laplace-precision}--\eqref{eq:laplace-covariance}.


\section*{3. The Two Corrections Under Misspecification: Overview}

The Laplace approximation \eqref{eq:laplace-precision} decomposes the
log-evidence as a sum of three contributions: the \emph{peak term}
$\log L(\hat\theta) + \log\pi(\hat\theta)$, an additive constant
$\tfrac{p}{2}\log(2\pi)$, and the \emph{volume term}
$-\tfrac{1}{2}\log\det(\boldsymbol\Lambda) =
\tfrac{1}{2}\log\det(\boldsymbol\Sigma)$. Under correct likelihood
specification at the population level, all three are evaluated consistently:
$\mathbf{JT}_{\mathrm{lik}} = \mathbf H_{\mathrm{lik}}$ (the Fisher identity)
propagates additively to $\mathbf C_{\mathrm{post}} = \mathbf I$
(Section~4.1 of the main supplement), so the naive Bayesian Laplace evidence
$\log Z_{\mathrm{F}}$ defined below is a valid asymptotic approximation.

Under likelihood misspecification, by contrast,
$\mathbf{JT}_{\mathrm{lik}} \neq \mathbf H_{\mathrm{lik}}$ and
$\mathbf C_{\mathrm{post}} \neq \mathbf I$. \emph{Both} the peak and the
volume term of the naive evidence are then biased, and each requires its
own correction. The two corrections are algebraically independent and
extract complementary scalar summaries of the same eigenvalue spectrum
$\mathrm{eig}(\mathbf C_{\mathrm{post}})$:
\begin{itemize}
\item The \textbf{peak correction} (Section~4 below) re-weights the
contribution $\log L(\hat\theta)$ via a learning rate
$\alpha_\star^{\mathrm{post}} = p / \mathrm{tr}(\mathbf C_{\mathrm{post}})$,
the \emph{arithmetic} mean of the eigenvalues of $\mathbf C_{\mathrm{post}}$.
Its leading-order effect on $\log Z$ is the additive term
$(\alpha_\star^{\mathrm{post}}-1)\log L(\hat\theta)$, which is $O(T)$ under
any persistent misspecification.
\item The \textbf{volume correction} (Sections~5--6 below) shifts the volume
term by $\tfrac{1}{2}\log\det(\mathbf C_{\mathrm{post}})$, the \emph{geometric}
mean of the eigenvalues. This shift is $O(1)$ in $T$ and survives even when
the peak correction is null (e.g.\ anti-isotropic distortions with
$\det \mathbf C_{\mathrm{post}} = 1$ but $\mathrm{tr}\mathbf C_{\mathrm{post}}
\neq p$).
\end{itemize}
The combined evidence $\log Z_{\mathrm{comb}}$ of Section~7 below
incorporates both corrections; the evidence-interval framework of
Section~10 reports $\log Z_{\mathrm{F}}$, $\log Z_{\mathrm{J}}$,
$\log Z_{\mathrm{DCC}}$, and $\log Z_{\mathrm{comb}}$ jointly as a
calibration-uncertainty band.

The reordering of this supplement relative to standard treatments
(\texttt{supplement\_evidence\_correction.tex} in the Likelihood folder
develops only the volume correction, and as a self-contained piece;
Bissiri et~al.\ 2016 develop the peak correction without the
sandwich-covariance reading) is deliberate: in our framework, the peak
correction is the \emph{primary} misspecification adjustment because it
acts on the likelihood value $L(\hat\theta)$ itself, while the volume
correction is a secondary adjustment on the derived posterior covariance.
Sections~4--7 develop the two corrections in this logical order; Sections
8--12 then characterise the resulting evidence and its diagnostic
properties.


\section*{4. Peak-Term Correction via Generalized Bayes Calibration}

Sections~5--6 (volume correction) and Section~7 (combined evidence) make sense only after the peak term has been calibrated, so we develop the peak correction first. Section~3 above gave a brief overview of the dual correction structure; this section derives the peak correction from the generalized-Bayes / power-posterior machinery of Bissiri, Holmes, and Walker (2016) and calibrates the tempering exponent $\alpha$ against $\mathbf C_{\mathrm{post}}$ via trace matching. The companion volume correction $\tfrac{1}{2}\log\det(\mathbf C_{\mathrm{post}})$ developed in Sections~5--6 uses the geometric-mean summary of the eigenvalues of $\mathbf C_{\mathrm{post}}$; the peak correction here uses the arithmetic-mean summary. The two are algebraically independent readouts of the same spectrum (see Section~4.6 below).

\subsection*{4.1 Motivation: peak vs.\ volume revisited}

The peak--volume decomposition \eqref{eq:peak-volume-split} of Section~11.3
makes the structural asymmetry of \eqref{eq:main-result} explicit: the IDM
correction reshapes the volume term only, while the $O(T)$ peak term is
unmodified. Under correct specification this is the right answer
($\mathbf C_{\mathrm{post}} = \mathbf I$ by Section~4.1 of the main supplement;
no volume correction needed either; see Section~8.3). Under misspecification with
$\mathbf{JT}_{\mathrm{lik}} \succ \mathbf H_{\mathrm{lik}}$, however, the
likelihood is effectively over-counted at $\hat\theta$: it acts as though it had
fewer independent observations than $T$, and the peak's contribution to $\log Z$
should be re-weighted accordingly. The volume correction alone cannot diagnose
this---it reads zero on any spectrum of $\mathbf C_{\mathrm{post}}$ with
$\det = 1$, including the anti-isotropic example $\mathbf C_{\mathrm{post}}
= \mathrm{diag}(2,\tfrac{1}{2})$ already flagged in Section~8.1 as motivation
for trace-based diagnostics, and revisited in \S10.6 below.

\subsection*{4.2 Generalized Bayes posterior and the power-likelihood}

The generalized-Bayes update of Bissiri, Holmes, and Walker (2016) replaces
$L(\theta)$ in the posterior by $L(\theta)^\alpha$ with $\alpha > 0$ a learning
rate / calibration scalar:
\begin{equation}
\pi_\alpha(\theta\mid D)
\;\propto\;
L(\theta)^\alpha\,\pi(\theta).
\label{eq:power-posterior}
\end{equation}
Under their coherence axioms (and the additivity of the loss across data), this
power-posterior is the unique coherent Bayesian update when the loss $-\log L$
is believed only up to a known scale; classical Bayes is the $\alpha = 1$
special case. The same calibrated-evidence object underlies the loss-likelihood
bootstrap of Lyddon, Holmes, and Walker (2019) and, in the information-criterion
literature, TIC (Takeuchi, 1976) and QAIC (Burnham and Anderson, 2002). See
Varin, Reid, and Firth (2011) for the composite-likelihood ancestor.

\subsection*{4.3 $\alpha$-tilted posterior Hessian, score, and Laplace evidence}

Differentiating $-\log\pi_\alpha(\theta\mid D) = -\alpha\log L(\theta)
- \log\pi(\theta)$ with respect to $\theta$, the prior score
$s_{\pi} = \nabla\log\pi$ is unchanged while the likelihood score $s_{\mathrm{lik}}
= \nabla\log L$ is multiplied by $\alpha$. The one-line variance calculation is
\[
\nabla(-\log\pi_\alpha) = -\alpha\, s_{\mathrm{lik}} - s_\pi,
\qquad
\mathrm{Var}_D\bigl(\nabla(-\log\pi_\alpha)\bigr) = \alpha^2\,\mathrm{Var}_D(s_{\mathrm{lik}})
= \alpha^2\,\mathbf{JT}_{\mathrm{lik}},
\]
since the prior gradient is deterministic given $\theta$. Under the additive
convention of \eqref{eq:setup-decomp}---which \emph{defines} $\mathbf{JT}_{\mathrm{post}}
:= \mathbf{JT}_{\mathrm{lik}} + \mathbf H_{\mathrm{prior}}$ (a definitional
convention, not a derived variance, because $\mathbf H_{\mathrm{prior}}$
contributes zero score variance)---the $\alpha$-tilted analogue is similarly
\emph{defined} by
\begin{align}
\text{Hessian of } -\log\pi_\alpha
&\;=\;
\alpha\,\mathbf H_{\mathrm{lik}} + \mathbf H_{\mathrm{prior}}
\;=:\;
\mathbf H_{\mathrm{post},\alpha},
\label{eq:Hpost-alpha}
\\[2pt]
\mathbf{JT}_{\mathrm{post},\alpha}
&\;:=\;
\alpha^2\,\mathbf{JT}_{\mathrm{lik}} + \mathbf H_{\mathrm{prior}},
\label{eq:JTpost-alpha}
\end{align}
so that $\alpha = 1$ recovers $\mathbf H_{\mathrm{post}}$ and
$\mathbf{JT}_{\mathrm{post}}$ of \eqref{eq:setup-decomp}. (We use the
comma-separated subscript $M_{a,\alpha}$ throughout \S10 to denote the
$\alpha$-tilted version of object $M_a$.) Setting
$\hat\theta_\alpha := \arg\max_\theta [\alpha\log L(\theta) + \log\pi(\theta)]$
and Laplace-expanding around it with the sandwich covariance
$\Sigma_{\mathrm{DCC},\mathrm{post},\alpha}
= \mathbf H_{\mathrm{post},\alpha}^{-1}\,\mathbf{JT}_{\mathrm{post},\alpha}\,
\mathbf H_{\mathrm{post},\alpha}^{-1}$ in place of $\mathbf H_{\mathrm{post}}^{-1}$
yields the $\alpha$-tilted DCC evidence
\begin{equation}
\log Z_\alpha
\;=\;
\alpha\,\log L(\hat\theta_\alpha)
\;+\;
\log\pi(\hat\theta_\alpha)
\;+\;
\tfrac{p}{2}\log(2\pi)
\;-\;
\log\det\!\bigl(\alpha\mathbf H_{\mathrm{lik}} + \mathbf H_{\mathrm{prior}}\bigr)
\;+\;
\tfrac{1}{2}\log\det\!\bigl(\alpha^2\mathbf{JT}_{\mathrm{lik}} + \mathbf H_{\mathrm{prior}}\bigr).
\label{eq:logZ-alpha}
\end{equation}
This is the \emph{sandwich} Laplace form (not the naive $-\tfrac{1}{2}\log\det
\mathbf H_{\mathrm{post},\alpha}$): the half-power in front of the log-det of
the sandwich, $\tfrac{1}{2}\log\det\Sigma_{\mathrm{DCC},\mathrm{post},\alpha}
= -\log\det \mathbf H_{\mathrm{post},\alpha} + \tfrac{1}{2}\log\det
\mathbf{JT}_{\mathrm{post},\alpha}$, accounts for the apparent factor-of-2
discrepancy with the naive Laplace, exactly as in \eqref{eq:SigmaDCCpost} and
\eqref{eq:DCC-recall} of Sections~1 and~3.2. At $\alpha = 1$,
\eqref{eq:logZ-alpha} collapses to the DCC evidence $\log Z_{\mathrm{DCC}}$ of
Sections~3.2 and~7.1.

\subsection*{4.4 Posterior trace-matching calibration $\alpha_\star^{\mathrm{post}} = p / \mathrm{tr}(\mathbf C_{\mathrm{post}})$}

We calibrate $\alpha$ by matching the $\alpha$-tilted effective parameter count
against the nominal dimension $p$. The trace of the $\alpha$-tilted information
distortion is the canonical TIC-style summary of the effective parameter count
(Takeuchi, 1976; Burnham and Anderson, 2002; Lyddon, Holmes, and Walker, 2019,
\S3): trace-matching is one of several conceivable calibrations (det-matching
or bootstrap-based alternatives would yield different $\alpha_\star$), and we
adopt it here for direct continuity with the established information-criterion
literature.

Define the $\alpha$-tilted posterior information distortion in direct analogy
with \eqref{eq:Cpost-def},
\begin{equation}
\mathbf C_{\mathrm{post},\alpha}
\;:=\;
\mathbf H_{\mathrm{post},\alpha}^{-1/2}\,
\mathbf{JT}_{\mathrm{post},\alpha}\,
\mathbf H_{\mathrm{post},\alpha}^{-1/2}
\;=\;
(\alpha\mathbf H_{\mathrm{lik}} + \mathbf H_{\mathrm{prior}})^{-1/2}\,
(\alpha^2\mathbf{JT}_{\mathrm{lik}} + \mathbf H_{\mathrm{prior}})\,
(\alpha\mathbf H_{\mathrm{lik}} + \mathbf H_{\mathrm{prior}})^{-1/2}.
\label{eq:Cpost-alpha}
\end{equation}
The exact trace-matching condition is the implicit nonlinear equation
\begin{equation}
\mathrm{tr}\!\bigl(\mathbf C_{\mathrm{post},\alpha}\bigr) \;=\; p.
\label{eq:trace-match-exact}
\end{equation}
Two simplifying observations:
\begin{itemize}
\item Under correct specification ($\mathbf{JT}_{\mathrm{lik}}
= \mathbf H_{\mathrm{lik}}$ propagating additively to $\mathbf{JT}_{\mathrm{post}}
= \mathbf H_{\mathrm{post}}$; Section~3.1 of the main supplement), substituting
$\alpha = 1$ into \eqref{eq:Cpost-alpha} gives
$\mathbf C_{\mathrm{post},1} = \mathbf H_{\mathrm{post}}^{-1/2}
\mathbf H_{\mathrm{post}} \mathbf H_{\mathrm{post}}^{-1/2} = \mathbf I$
and $\mathrm{tr}(\mathbf C_{\mathrm{post},1}) = p$, so $\alpha = 1$ is an
\emph{exact} solution of \eqref{eq:trace-match-exact} for any proper Gaussian
prior. This preserves the marquee correct-specification property cleanly.
\item In the data-dominated limit
$\mathbf H_{\mathrm{prior}} \ll \alpha\mathbf H_{\mathrm{lik}}$ (and
$\mathbf H_{\mathrm{prior}} \ll \alpha^2\mathbf{JT}_{\mathrm{lik}}$),
\eqref{eq:Cpost-alpha} reduces at leading order to
\begin{equation}
\mathbf C_{\mathrm{post},\alpha}
\;\approx\;
(\alpha\mathbf H_{\mathrm{lik}})^{-1/2}(\alpha^2\mathbf{JT}_{\mathrm{lik}})
(\alpha\mathbf H_{\mathrm{lik}})^{-1/2}
\;=\;
\alpha\,\mathbf C_{\mathrm{lik}},
\label{eq:Cpost-alpha-asymptote}
\end{equation}
i.e.\ the limit object is the \emph{likelihood-side} distortion
$\mathbf C_{\mathrm{lik}}$, not $\mathbf C_{\mathrm{post}}$ (the two coincide
in this limit, $\mathbf C_{\mathrm{post}} \to \mathbf C_{\mathrm{lik}}$).
\end{itemize}
The leading-order solution of \eqref{eq:trace-match-exact} in the data-dominated
limit, using either $\mathrm{tr}(\mathbf C_{\mathrm{post},\alpha})
\approx \alpha\,\mathrm{tr}(\mathbf C_{\mathrm{lik}}) \approx \alpha\,
\mathrm{tr}(\mathbf C_{\mathrm{post}})$, is the boxed posterior calibration
\begin{equation}
\boxed{\;
\alpha_\star^{\mathrm{post}}
\;:=\;
\frac{p}{\mathrm{tr}(\mathbf C_{\mathrm{post}})}.
\;}
\label{eq:alpha-star-post}
\end{equation}
We use the explicit ``post'' superscript to mark this as the
posterior-supplement calibration; the likelihood-folder counterpart
$\alpha_\star^{\mathrm{lik}} = p/\mathrm{tr}(\mathbf C_{\mathrm{lik}})$ (\S10.9)
uses the same boxed identity with the prior set to zero. The defining
properties of \eqref{eq:alpha-star-post} mirror those of
$\log\det(\mathbf C_{\mathrm{post}})$ in Sections~5.1--5.4:
\begin{itemize}
\item \emph{Correct specification.} As noted above,
$\mathbf C_{\mathrm{post},1} = \mathbf I$ \emph{exactly} at $\alpha = 1$ for
any proper Gaussian prior; the leading-order formula
\eqref{eq:alpha-star-post} also returns $\alpha_\star^{\mathrm{post}} = 1$
since $\mathrm{tr}(\mathbf C_{\mathrm{post}}) = p$. This is the trace analogue
of the prior-cancellation property of \S5.3.
\item \emph{Variance inflation.} When
$\mathbf{JT}_{\mathrm{lik}} \succ \mathbf H_{\mathrm{lik}}$,
$\mathrm{tr}(\mathbf C_{\mathrm{post}}) > p$ and
$\alpha_\star^{\mathrm{post}} < 1$: the power-likelihood weight is reduced, so
each observation enters the calibrated evidence at less than its face value
(an \emph{effective-sample-size} reduction, not necessarily a reduction of
$\log Z$; see \S10.5 on sign-dependence).
\item \emph{Deflation.} When
$\mathbf{JT}_{\mathrm{lik}} \prec \mathbf H_{\mathrm{lik}}$,
$\mathrm{tr}(\mathbf C_{\mathrm{post}}) < p$ and
$\alpha_\star^{\mathrm{post}} > 1$: the power-likelihood weight is increased
(observations are over-weighted relative to face value).
\item \emph{Well-posedness, not informativity.}
$\alpha_\star^{\mathrm{post}}$ is mathematically well-defined whenever
$\mathrm{tr}(\mathbf C_{\mathrm{post}}) > 0$, i.e.\ whenever
$\mathbf H_{\mathrm{post}} \succ 0$, i.e.\ under the dominance condition of
Section~2.1 of the main supplement
(label \texttt{eq:dominance-condition}). No subspace projection or
retained-eigenspace machinery is required, in contrast to the likelihood-only
version (see \S10.9). Well-posedness does not imply informativity: under a
sufficiently strong prior, $\mathbf C_{\mathrm{post}} \to \mathbf I$ even
under misspecification (\S10.4.1 below).
\end{itemize}

\subsubsection*{10.4.1 Strong-prior boundary: $\alpha_\star^{\mathrm{post}} \to 1$ does not imply correct spec}

The same residual bound used in Section~8.4 to flag the strong-prior boundary
of the volume correction applies verbatim to the trace summary. By Section~4.2
of the main supplement (label
\texttt{eq:posterior-c-residual-bound}),
\[
\|\mathbf C_{\mathrm{post}} - \mathbf I\|
\;\le\;
\lambda_{\min}(\mathbf H_{\mathrm{post}})^{-1}\,
\|\Delta\mathbf{H}^{\mathrm{lik}}_{\mathrm{resid}}\|
\;\xrightarrow[\lambda_{\min}(\mathbf H_{\mathrm{prior}})\to\infty]{}\;
0,
\]
which drives both $\det(\mathbf C_{\mathrm{post}}) \to 1$ and
$\mathrm{tr}(\mathbf C_{\mathrm{post}}) \to p$. Hence $\alpha_\star^{\mathrm{post}}
\to 1$ in the strong-prior limit \emph{even under arbitrary misspecification},
exactly paralleling the $\Delta\log Z \to 0$ ambiguity of \S5.4. The diagnostic
$\alpha_\star^{\mathrm{post}} \approx 1$ alone is therefore insufficient to
certify correct specification; resolving the ambiguity requires varying the
prior strength or inspecting per-eigenvalue diagnostics of
$\mathbf C_{\mathrm{lik}}$ on a retained subspace, exactly as in \S5.4 for the
volume correction. See \S10.10 for operational guidance.

\subsection*{4.5 The combined peak$+$volume evidence}

Inserting $\alpha_\star^{\mathrm{post}}$ from \eqref{eq:alpha-star-post} into
\eqref{eq:logZ-alpha} and using \eqref{eq:main-result} for the volume piece, we
obtain the combined-corrected posterior evidence
\begin{equation}
\log Z_{\mathrm{comb}}
\;=\;
\alpha_\star^{\mathrm{post}}\log L(\hat\theta_{\alpha_\star^{\mathrm{post}}})
\;+\;
\log\pi(\hat\theta_{\alpha_\star^{\mathrm{post}}})
\;+\;
\tfrac{p}{2}\log(2\pi)
\;-\;
\log\det\!\bigl(\alpha_\star^{\mathrm{post}}\mathbf H_{\mathrm{lik}}
+ \mathbf H_{\mathrm{prior}}\bigr)
\;+\;
\tfrac{1}{2}\log\det\!\bigl((\alpha_\star^{\mathrm{post}})^2
\mathbf{JT}_{\mathrm{lik}} + \mathbf H_{\mathrm{prior}}\bigr).
\label{eq:logZ-comb}
\end{equation}
Two reductions to leading order:
\begin{itemize}
\item \emph{MAP shift.} In the data-dominated limit, the gradient of
$\alpha\log L + \log\pi$ at $\hat\theta_\alpha$ vanishes while that of
$\log L + \log\pi$ at $\hat\theta$ vanishes; expanding,
$\hat\theta_{\alpha_\star^{\mathrm{post}}} - \hat\theta
= (\alpha_\star^{\mathrm{post}} - 1)\,
\mathbf H_{\mathrm{post}}^{-1}\,s_\pi(\hat\theta) + \cdots$. Since
$s_\pi(\hat\theta) = O(1)$ while $\mathbf H_{\mathrm{post}} = O(T)$, the MAP
shift is $O((\alpha_\star^{\mathrm{post}}-1)/T)$; the induced peak shift
$\log L(\hat\theta_{\alpha_\star^{\mathrm{post}}}) - \log L(\hat\theta)$ is
$O((\alpha_\star^{\mathrm{post}}-1)/T) \cdot s_{\mathrm{lik}}(\hat\theta)
+ O((\alpha_\star^{\mathrm{post}}-1)^2/T)$, and $s_{\mathrm{lik}}(\hat\theta)
= -s_\pi(\hat\theta) = O(1)$ at the classical MAP, so the peak shift is
$O((\alpha_\star^{\mathrm{post}}-1)/T) = o(1)$. We may therefore substitute
$\log L(\hat\theta_{\alpha_\star^{\mathrm{post}}}) \to \log L(\hat\theta)$ and
$\log\pi(\hat\theta_{\alpha_\star^{\mathrm{post}}}) \to \log\pi(\hat\theta)$ in
\eqref{eq:logZ-comb} with errors that are negligible against the leading
$(\alpha_\star^{\mathrm{post}}-1)\log L(\hat\theta) = O(T)$ contribution.
\item \emph{Prior-in-log-det cancellation.} In the same data-dominated limit
the prior terms factor out of both log-dets,
$\log\det(\alpha\mathbf H_{\mathrm{lik}} + \mathbf H_{\mathrm{prior}})
= p\log\alpha + \log\det\mathbf H_{\mathrm{lik}}
+ \mathrm{tr}\bigl[(\alpha\mathbf H_{\mathrm{lik}})^{-1}\mathbf H_{\mathrm{prior}}\bigr]
+ \cdots$ and similarly for the $\mathbf{JT}$ log-det, so the
$\alpha_\star^{\mathrm{post}}$-dependent log-det contributions cancel at
leading order:
\begin{equation}
-\,p\log\alpha_\star^{\mathrm{post}}
\;+\;
\tfrac{1}{2}\bigl(2p\log\alpha_\star^{\mathrm{post}}\bigr)
\;=\;
0.
\label{eq:logdet-cancel}
\end{equation}
\end{itemize}
Combining the two reductions, \eqref{eq:logZ-comb} reduces to
\begin{equation}
\log Z_{\mathrm{comb}}
\;\approx\;
\log Z_{\mathrm{DCC}}
\;+\;
(\alpha_\star^{\mathrm{post}} - 1)\,\log L(\hat\theta)
\;+\;
O\!\bigl(\mathrm{tr}[\mathbf H_{\mathrm{lik}}^{-1}\mathbf H_{\mathrm{prior}}]/
\alpha_\star^{\mathrm{post}}\bigr).
\label{eq:logZ-comb-leading}
\end{equation}
The combined correction is therefore the IDM volume shift of
\eqref{eq:main-result} plus a peak re-weight of size
$(\alpha_\star^{\mathrm{post}} - 1)\log L(\hat\theta)$. Because
$\log L(\hat\theta) = O(T)$ while
$\alpha_\star^{\mathrm{post}} - 1 = O(1)$ in $T$ (it is a function of
intensive moments of $\mathbf C_{\mathrm{post}}$), the peak correction is
$O(T)$ in general and dominates the $O(1)$ volume correction in the
data-dominated misspecified regime. This is the structural reordering flagged
in the chapeau: under stationary misspecification the spread
$\log Z_{\mathrm{comb}} - \log Z_{\mathrm{DCC}}$ scales with $T$, in contrast
with the $O(1)$ spread of $(\log Z_F, \log Z_J, \log Z_{\mathrm{DCC}})$
analysed in \S7.2 (which is also $O(1)$ via \eqref{eq:correction-O1}). The
$O(1)$ statements of \S\S8.3--8.4 thus apply to the volume-only correction;
they are not contradicted by \eqref{eq:logZ-comb-leading} but supplemented by
it.

\subsubsection*{10.5.1 Sign of the peak shift on $\log Z$}

The directional language \emph{discounted}/\emph{boosted} applies to the
\emph{weight} placed on the log-likelihood by the power-posterior, not
directly to the sign of the shift on $\log Z$. Whether
$(\alpha_\star^{\mathrm{post}} - 1)\log L(\hat\theta)$ raises or lowers
$\log Z_{\mathrm{comb}}$ relative to $\log Z_{\mathrm{DCC}}$ depends on the
sign of $\log L(\hat\theta)$, which is parameterisation-dependent (densities
above one give $\log L > 0$, probabilities below one give $\log L < 0$). The
unambiguous reading is in effective-sample-size terms:
$\alpha_\star^{\mathrm{post}} < 1$ means the data are treated as carrying
fewer effective independent observations than $T$, regardless of how this
translates to a $\log Z$ shift. For Bayes-factor comparisons the
decision-relevant quantity is the \emph{difference}
$(\alpha_\star^{\mathrm{post}, A} \log L_A(\hat\theta_A))
- (\alpha_\star^{\mathrm{post}, B} \log L_B(\hat\theta_B))$, where the
two calibrations are different and the parameterisation-dependent sign
issue is partly absorbed.

\subsection*{4.6 Two algebraically independent moments of the same spectrum}

The volume correction of Sections~5--6 and~10 and the peak correction of \S10.4 use
\emph{different scalar summaries of the same eigenvalue spectrum} of
$\mathbf C_{\mathrm{post}}$. Section~8.1 already noted that the determinant
readout is blind to anti-isotropic distortions such as
$\mathrm{diag}(2,\tfrac{1}{2})$ and recommended trace-based diagnostics; the
calibration \eqref{eq:alpha-star-post} is precisely such a diagnostic, making
explicit which scalar summary of $\mathrm{eig}(\mathbf C_{\mathrm{post}})$ the
peak correction probes. Let $\lambda_1,\ldots,\lambda_p > 0$ be those
eigenvalues. Then
\begin{align}
\text{IDM volume correction}
&\;=\;
\tfrac{1}{2}\log\det(\mathbf C_{\mathrm{post}})
\;=\;
\tfrac{p}{2}\,\log\overline\lambda_{\mathrm{geom}},
&
\overline\lambda_{\mathrm{geom}}
&\;=\;
\Bigl(\prod_i\lambda_i\Bigr)^{1/p},
\\[2pt]
\text{$\alpha$-calibration scalar}
&\;=\;
\alpha_\star^{\mathrm{post}}
\;=\;
\frac{1}{\overline\lambda_{\mathrm{arith}}},
&
\overline\lambda_{\mathrm{arith}}
&\;=\;
\frac{1}{p}\sum_i\lambda_i
\;=\;
\frac{\mathrm{tr}(\mathbf C_{\mathrm{post}})}{p}.
\end{align}
The IDM correction probes the geometric mean of the spectrum, the peak
correction the arithmetic mean. By the AM--GM inequality
$\overline\lambda_{\mathrm{arith}} \ge \overline\lambda_{\mathrm{geom}}$, with
equality iff $\mathbf C_{\mathrm{post}} = c\mathbf I$. The two corrections are
therefore not double-counting: they are algebraically independent symmetric
functions of the spectrum (we avoid the word ``orthogonal,'' which has no
Hilbert-space meaning here---the relation is the AM--GM inequality, not an
inner-product orthogonality).

\emph{Diagonal example.} Take $\mathbf C_{\mathrm{post}}
= \mathrm{diag}(2,\tfrac{1}{2})$ with $p = 2$:
\[
\det(\mathbf C_{\mathrm{post}}) \;=\; 1
\;\Rightarrow\;
\tfrac{1}{2}\log\det(\mathbf C_{\mathrm{post}}) \;=\; 0
\qquad\text{(IDM volume correction vanishes)},
\]
\[
\mathrm{tr}(\mathbf C_{\mathrm{post}}) \;=\; 2.5
\;\Rightarrow\;
\alpha_\star^{\mathrm{post}} \;=\; \frac{2}{2.5} \;=\; 0.8
\qquad\text{(peak weight reduced by 20\%)}.
\]
This anti-isotropic misspecification---one direction inflated, another
deflated---is invisible to the determinant readout of \eqref{eq:main-result}
and to the evidence-interval width
$W_Z = |\log\det(\mathbf C_{\mathrm{post}})|$ of \S7, but is detected by the
trace-based $\alpha_\star^{\mathrm{post}}$. Reporting both moments is strictly
more informative than reporting either alone.

\subsection*{4.7 Companion readout in the information-gain supplement}

The peak--volume duality has a same-folder companion in
\texttt{supplement\_information\_gain.tex}: the gap
$\Delta_{\mathbf G} = -\log\det(\mathbf C_{\mathrm{post}})$ captured there
(see Eq.~\eqref{eq:DeltaG-identity} of that supplement; sign convention as
stated in the companion's Section~7) is the geometric-mean readout of the
same spectrum, and the calibration scalar
$\alpha_\star^{\mathrm{post}} = p/\mathrm{tr}(\mathbf C_{\mathrm{post}})$ is
its arithmetic-mean companion. Together they form a two-axis diagnostic
$(\overline\lambda_{\mathrm{arith}},\overline\lambda_{\mathrm{geom}})$ in
which correct specification corresponds to $(1,1)$, isotropic inflation to a
movement along the diagonal, and anti-isotropic distortion to a movement
away from the diagonal. We recommend reporting both axes jointly whenever
$\mathbf C_{\mathrm{post}}$ is computed.

\subsection*{4.8 Composition with the evidence-interval framework}

Section~7 introduced the evidence interval
$(\log Z_F,\,\log Z_J,\,\log Z_{\mathrm{DCC}})$ summarising determinant-level
calibration uncertainty. The combined estimator $\log Z_{\mathrm{comb}}$ of
\eqref{eq:logZ-comb} augments this to a four-point set
\begin{equation}
(\log Z_F,\;\log Z_J,\;\log Z_{\mathrm{DCC}},\;\log Z_{\mathrm{comb}}),
\label{eq:four-point-set}
\end{equation}
whose members partition the correction structure as follows:
\begin{center}
\renewcommand{\arraystretch}{1.2}
\begin{tabular}{lccc}
\hline
estimator & volume correction & peak correction & calibration scalar \\
\hline
$\log Z_F$           & none                                   & none & --- \\
$\log Z_J$           & $\mathbf{JT}_{\mathrm{post}}$-only     & none & --- \\
$\log Z_{\mathrm{DCC}}$ & $\tfrac{1}{2}\log\det(\mathbf C_{\mathrm{post}})$ & none & geometric-mean only \\
$\log Z_{\mathrm{comb}}$ & full sandwich, $\alpha$-tilted    & $(\alpha_\star^{\mathrm{post}}-1)\log L(\hat\theta)$ & geometric $+$ arithmetic \\
\hline
\end{tabular}
\end{center}
The pairwise differences among $(\log Z_F,\log Z_J,\log Z_{\mathrm{DCC}})$ are
all $O(1)$ (\S7.2 and \eqref{eq:correction-O1}); the spread
$\log Z_{\mathrm{comb}} - \log Z_{\mathrm{DCC}}$ is, by
\eqref{eq:logZ-comb-leading}, the $O(T)$ peak shift
$(\alpha_\star^{\mathrm{post}} - 1)\log L(\hat\theta)$. The expanded set
\eqref{eq:four-point-set} is therefore not four members of a single $O(1)$
interval but a two-tier structure: an $O(1)$ inner interval
$(\log Z_F, \log Z_J, \log Z_{\mathrm{DCC}})$ summarising determinant-level
calibration, plus an outer $\log Z_{\mathrm{comb}}$ that incorporates the
$O(T)$ peak re-weight. For sensitivity analysis the four-point set carries
strictly more information than any of its members; as a single headline
quantity it should not be collapsed.

A subsidiary remark on \emph{within-axis} effective sample sizes. The
calibration uses $\mathbf{JT}_{\mathrm{post}}$ (the total score-fluctuation
budget) and not $\mathbf{JS}_{\mathrm{post}}$ (the per-sample/diagonal block
of Section~6.3 used to define $\mathbf R_{\mathrm{post}}$ and
$T_{\mathrm{eff}}^{\mathrm{post}}$). The $\mathbf C_{\mathrm{post}}$-based
$N_{\mathrm{eff}}^{Z}$ of \S6.3 governs the volume correction; the
$\mathbf C_{\mathrm{post}}$-based $\alpha_\star^{\mathrm{post}}$ governs the
peak correction; the $\mathbf R_{\mathrm{post}}$-based
$T_{\mathrm{eff}}^{\mathrm{post}}$ governs score-vs-Fisher concentration. These
are three distinct effective-sample-size objects with three distinct
operational meanings and should not be conflated.

\subsection*{4.9 Singular $\mathbf H_{\mathrm{lik}}$: why the posterior version is preferred}

The likelihood-only counterpart of the calibration \eqref{eq:alpha-star-post} is
$\alpha_\star^{\mathrm{lik}} = p / \mathrm{tr}(\mathbf C_{\mathrm{lik}})$ with
$\mathbf C_{\mathrm{lik}} = \mathbf H_{\mathrm{lik}}^{-1/2}
\mathbf{JT}_{\mathrm{lik}} \mathbf H_{\mathrm{lik}}^{-1/2}$ (the standard
textbook TIC of Takeuchi (1976) / QAIC of Burnham and Anderson (2002)). The
companion section in \texttt{supplement\_evidence\_correction.tex}
(the likelihood-folder counterpart of this section) develops the
likelihood-only version in its own notation.

\emph{Notation bridge.} The likelihood-folder supplement
\texttt{supplement\_evidence\_correction.tex} uses bare
$\mathbf H, \mathbf C, \mathbf J_{\mathrm{total}}$ where this document writes
$\mathbf H_{\mathrm{lik}}, \mathbf C_{\mathrm{lik}}, \mathbf{JT}_{\mathrm{lik}}$;
the identifications are
\[
\mathbf H \equiv \mathbf H_{\mathrm{lik}},
\quad
\mathbf C \equiv \mathbf C_{\mathrm{lik}},
\quad
\mathbf J_{\mathrm{total}} \equiv \mathbf{JT}_{\mathrm{lik}},
\]
and the calibration scalar in that folder is written
$\alpha_\star^{\mathrm{lik}}$ (with the explicit ``lik'' superscript when
folder-disambiguation is needed).

The likelihood form requires $\mathbf H_{\mathrm{lik}} \succ 0$, which fails
generically near the boundary of $\Theta$, in finite-difference Hessian
implementations with indefinite spectrum, and in any sloppy-model regime where
some eigenvalues of $\mathbf H_{\mathrm{lik}}$ are vanishingly small. The
singular-$\mathbf H_{\mathrm{lik}}$ subspace-projection machinery used by the
likelihood-folder supplement (its Setup section and its singular-Fisher
treatment) is therefore required there but \emph{not} here.

The posterior calibration \eqref{eq:alpha-star-post} is regularised by the
prior: it is well-defined whenever $\mathbf H_{\mathrm{post}} \succ 0$, i.e.\
under the dominance condition of Section~2.1 of the main supplement
(label \texttt{eq:dominance-condition}). In the data-dominated limit
$\mathbf H_{\mathrm{prior}} \ll \mathbf H_{\mathrm{lik}}$ the two coincide
($\mathbf C_{\mathrm{post}} \to \mathbf C_{\mathrm{lik}}$ and
$\alpha_\star^{\mathrm{post}} \to \alpha_\star^{\mathrm{lik}}$), so the
posterior form loses nothing in the well-behaved regime and remains operative
in the singular-$\mathbf H_{\mathrm{lik}}$ regime where the likelihood form is
undefined. The reader interested in the classical likelihood-only formulation
and its TIC/QAIC roots is referred to the corresponding section of
\texttt{supplement\_evidence\_correction.tex}.

\subsection*{4.10 When to apply, when to skip, when to worry}

Recommended practice:
\begin{itemize}
\item \emph{Apply} when $\alpha_\star^{\mathrm{post}} \lesssim 0.9$
(equivalently $\mathrm{tr}(\mathbf C_{\mathrm{post}})/p \gtrsim 1.1$) or
$\alpha_\star^{\mathrm{post}} \gtrsim 1.1$: the peak-term re-weight is then
comparable to or larger than the volume correction and changes the
leading-order $O(T)$ contribution to $\log Z$.
\item \emph{Skip} when $\alpha_\star^{\mathrm{post}} \approx 1$ \emph{and}
$|\log\det(\mathbf C_{\mathrm{post}})|$ small, \emph{and} the prior is not so
strong that both readouts are suppressed by the strong-prior mechanism of
\S10.4.1 (mirror of \S5.4). The disambiguation is the same as in \S5.4: vary
the prior strength. Both moments of the spectrum confirm
$\mathbf C_{\mathrm{post}} \approx \mathbf I$ only under a weak-prior reading;
under a strong prior they confirm only that the prior is dominant.
\item \emph{Worry} when the two moments disagree, e.g.\
$\alpha_\star^{\mathrm{post}}$ substantially different from $1$ but
$\log\det(\mathbf C_{\mathrm{post}}) \approx 0$ (the anti-isotropic case of
\S10.6), or vice versa: a single scalar correction is then insufficient and
per-eigenvalue diagnostics ($\mathrm{eig}(\mathbf C_{\mathrm{post}})$) should
be inspected.
\item \emph{Validity-band caveat for the alpha-Laplace.} The $\alpha$-tilted
Laplace expansion underlying \eqref{eq:logZ-alpha} is reliable while
$\alpha_\star^{\mathrm{post}}\mathbf H_{\mathrm{lik}}
+ \mathbf H_{\mathrm{prior}}$ is well-conditioned with respect to the
$\alpha$-tilted posterior's concentration around
$\hat\theta_{\alpha_\star^{\mathrm{post}}}$. Outside the band $0.5 \lesssim
\alpha_\star^{\mathrm{post}} \lesssim 2$ the leading-order expansion degrades:
$\alpha_\star^{\mathrm{post}} \to 0^+$ approaches the prior-only inference
limit (data effectively ignored, $L^{\alpha} \to 1$), and
$\alpha_\star^{\mathrm{post}} \gg 1$ amplifies any score--Fisher mismatch
beyond the regime where a single quadratic expansion is faithful. In the
extreme regimes, treat $\alpha_\star^{\mathrm{post}}$ as a flag rather than a
usable correction: severely small or severely large values signal that the
likelihood model itself should be revisited.
\item \emph{Caveat for Bayes-factor comparisons.} Because
$\alpha_\star^{\mathrm{post}}$ is data-driven, it can over-penalise complex
models (Gr\"unwald and van Ommen, 2017; M\"uller, 2013): a more flexible
model fits the noise, inflates $\mathbf{JT}_{\mathrm{lik}}$ relative to
$\mathbf H_{\mathrm{lik}}$, depresses $\alpha_\star^{\mathrm{post}}$, and also
raises $\log\det(\mathbf C_{\mathrm{post}})$. The two corrections then move
$\log Z_{\mathrm{comb}}$ in \emph{opposite directions}---the volume term
rewards the inflated determinant while the peak re-weight reduces the
log-likelihood weight---and the net effect is parameterisation-dependent
(it requires knowing the sign of $\log L(\hat\theta)$). The framing
``doubly penalised'' is therefore misleading; report both corrections rather
than a single net shift, and use $\log Z_{\mathrm{comb}}$ as a sensitivity
check alongside $\log Z_{\mathrm{DCC}}$, not as the headline quantity. The
full set \eqref{eq:four-point-set} carries strictly more information than any
of its members.
\item \emph{Joint diagnostic.} The minimal informative summary of the
distortion is the pair
$\bigl(\alpha_\star^{\mathrm{post}},\,
\tfrac{1}{2}\log\det(\mathbf C_{\mathrm{post}})\bigr)$:
arithmetic mean against geometric mean of the same spectrum, with the
strong-prior caveat of \S10.4.1 attached to both.
\end{itemize}

\paragraph{References for this section.}
Bissiri, P.~G., Holmes, C.~C., and Walker, S.~G.\ (2016) ``A general framework
for updating belief distributions,'' \emph{J.~R.~Statist.~Soc.~B} 78(5):
1103--1130.\;
Lyddon, S.~P., Holmes, C.~C., and Walker, S.~G.\ (2019) ``General Bayesian
updating and the loss-likelihood bootstrap,'' \emph{Biometrika} 106(2):
465--478.\;
Varin, C., Reid, N., and Firth, D.\ (2011) ``An overview of composite
likelihood methods,'' \emph{Statistica Sinica} 21: 5--42.\;
Gr\"unwald, P. and van Ommen, T.\ (2017) ``Inconsistency of Bayesian
inference for misspecified linear models,'' \emph{Bayesian Analysis} 12(4).\;
M\"uller, U.~K.\ (2013) ``Risk of Bayesian inference in misspecified
models,'' \emph{Econometrica} 81: 1805--1849.\;
Takeuchi, K.\ (1976) (TIC); Burnham, K.~P.\ and Anderson, D.~R.\ (2002)
\emph{Model Selection and Multimodel Inference} (QAIC).

\section*{5. Volume-Term Correction: Naive versus Corrected Laplace}

\subsection*{5.1 Naive Bayesian Laplace}

The standard Bayesian Laplace approximation identifies the local posterior precision
with the negative posterior Hessian at the MAP,
\begin{equation}
\boldsymbol\Lambda_{\mathrm{naive}}
\;=\;
\mathbf H_{\mathrm{post}},
\qquad
\Sigma_{\mathrm{naive}}
\;=\;
\mathbf H_{\mathrm{post}}^{-1},
\label{eq:naive-curvature}
\end{equation}
giving the naive log-evidence
\begin{equation}
\log Z_{\mathrm{naive}}
\;=\;
\log L(\hat\theta) + \log\pi(\hat\theta)
\;+\;
\tfrac{p}{2}\log(2\pi)
\;-\;
\tfrac{1}{2}\log\det(\mathbf H_{\mathrm{post}}).
\label{eq:naive-Z}
\end{equation}
This is the formula underlying BIC-style asymptotic approximations and the standard
Bayesian Laplace evidence.

\subsection*{5.2 Corrected Bayesian Laplace}

Under possible likelihood misspecification, the asymptotic covariance of the MAP estimator
is the sandwich (Section~6 of the main supplement)
\begin{equation}
\Sigma_{\mathrm{DCC},\mathrm{post}}
\;=\;
\mathbf H_{\mathrm{post}}^{-1}\,
\mathbf{JT}_{\mathrm{post}}\,
\mathbf H_{\mathrm{post}}^{-1},
\label{eq:DCC-recall}
\end{equation}
not $\mathbf H_{\mathrm{post}}^{-1}$. Replacing $\Sigma_{\mathrm{naive}}$ with
$\Sigma_{\mathrm{DCC},\mathrm{post}}$ in \eqref{eq:laplace-covariance} gives
\begin{equation}
\log Z_{\mathrm{corrected}}
\;=\;
\log L(\hat\theta) + \log\pi(\hat\theta)
\;+\;
\tfrac{p}{2}\log(2\pi)
\;+\;
\tfrac{1}{2}\log\det\!\left(
\mathbf H_{\mathrm{post}}^{-1}\,
\mathbf{JT}_{\mathrm{post}}\,
\mathbf H_{\mathrm{post}}^{-1}
\right).
\label{eq:corrected-Z}
\end{equation}
By multiplicativity of the determinant, $\det(\mathbf A\mathbf B\mathbf A) = \det(\mathbf A)^2 \det(\mathbf B)$ for any square $\mathbf A, \mathbf B$, hence
\begin{equation}
\log\det\!\left(
\mathbf H_{\mathrm{post}}^{-1}\,
\mathbf{JT}_{\mathrm{post}}\,
\mathbf H_{\mathrm{post}}^{-1}
\right)
\;=\;
\log\det(\mathbf{JT}_{\mathrm{post}})
\;-\;
2\log\det(\mathbf H_{\mathrm{post}}).
\label{eq:logdet-expansion}
\end{equation}

\section*{6. Volume-Term Correction: The Shift in Log-Evidence}

Subtract \eqref{eq:naive-Z} from \eqref{eq:corrected-Z} and use
\eqref{eq:logdet-expansion}:
\begin{align}
\Delta\log Z
\;&:=\;
\log Z_{\mathrm{corrected}} \;-\; \log Z_{\mathrm{naive}}
\nonumber\\[4pt]
&=\;
\tfrac{1}{2}\bigl[
\log\det(\mathbf{JT}_{\mathrm{post}})
\;-\;
2\log\det(\mathbf H_{\mathrm{post}})
\bigr]
\;-\;
\bigl[
-\tfrac{1}{2}\log\det(\mathbf H_{\mathrm{post}})
\bigr]
\nonumber\\[4pt]
&=\;
\tfrac{1}{2}\log\det(\mathbf{JT}_{\mathrm{post}})
\;-\;
\tfrac{1}{2}\log\det(\mathbf H_{\mathrm{post}})
\nonumber\\[4pt]
&=\;
\tfrac{1}{2}\log
\frac{\det(\mathbf{JT}_{\mathrm{post}})}
     {\det(\mathbf H_{\mathrm{post}})}.
\label{eq:DeltaZ-ratio}
\end{align}
Recognising the right-hand side as the log-determinant of the Posterior Information
Distortion,
\begin{equation}
\det(\mathbf C_{\mathrm{post}})
\;=\;
\det\!\left(
\mathbf H_{\mathrm{post}}^{-1/2}\,
\mathbf{JT}_{\mathrm{post}}\,
\mathbf H_{\mathrm{post}}^{-1/2}
\right)
\;=\;
\frac{\det(\mathbf{JT}_{\mathrm{post}})}
     {\det(\mathbf H_{\mathrm{post}})},
\label{eq:detCpost-identity}
\end{equation}
we obtain the central result of this supplement:
\begin{equation}
\boxed{\;
\Delta\log Z
\;=\;
\tfrac{1}{2}\log\det(\mathbf C_{\mathrm{post}}).
\;}
\label{eq:main-result}
\end{equation}

\paragraph{Direct parallel to the likelihood case.}
Result~\eqref{eq:main-result} is the direct posterior counterpart of the central result
of \texttt{supplement\_evidence\_correction.tex}, with the likelihood-side objects
$(\mathbf H, \mathbf J_{\mathrm{total}}, \mathbf C)$ replaced throughout by
$(\mathbf H_{\mathrm{post}}, \mathbf{JT}_{\mathrm{post}}, \mathbf C_{\mathrm{post}})$
of \eqref{eq:setup-decomp}--\eqref{eq:Cpost-def}. The proof is identical at the level of
algebra; the substantive content lies in what changes upon promotion from likelihood to
posterior, treated in the next section.


\section*{7. The Combined Evidence $\log Z_{\mathrm{comb}}$}

Sections~4 and 6 produced two independent corrections to the naive Bayesian
Laplace evidence: the peak correction
$(\alpha_\star^{\mathrm{post}}-1)\log L(\hat\theta) + O(p\log\alpha_\star^{\mathrm{post}})$
of Section~4.5 (data-dominated leading order) and the volume correction
$\tfrac{1}{2}\log\det(\mathbf C_{\mathrm{post}})$ of Section~6. The
\emph{combined evidence} is obtained by composing the two corrections
on the $\alpha$-tilted Laplace formula of Eq.~\eqref{eq:logZ-alpha}:
\begin{equation}
\boxed{\;
\log Z_{\mathrm{comb}}
\;:=\;
\alpha_\star^{\mathrm{post}}\,\log L(\hat\theta)
\;+\;
\log\pi(\hat\theta)
\;+\;
\tfrac{p}{2}\log(2\pi)
\;-\;
\log\det\!\bigl(\alpha_\star^{\mathrm{post}}\mathbf H_{\mathrm{lik}}
+ \mathbf H_{\mathrm{prior}}\bigr)
\;+\;
\tfrac{1}{2}\log\det\!\bigl(({\alpha_\star^{\mathrm{post}}})^2\,
\mathbf{JT}_{\mathrm{lik}} + \mathbf H_{\mathrm{prior}}\bigr).
\;}
\label{eq:logZ-comb-def}
\end{equation}
This is the formula obtained by substituting
$\alpha_\star^{\mathrm{post}} = p/\mathrm{tr}(\mathbf C_{\mathrm{post}})$
(Section~4.4) into \eqref{eq:logZ-alpha}. At $\alpha_\star^{\mathrm{post}} = 1$
(correct specification regime), $\log Z_{\mathrm{comb}}$ collapses to
$\log Z_{\mathrm{DCC}}$ of Section~5.2, and both reduce to the naive
$\log Z_{\mathrm{F}}$ when also $\mathbf C_{\mathrm{post}} = \mathbf I$.

\paragraph{Data-dominated leading order.}
In the regime $\mathbf H_{\mathrm{prior}} \ll \alpha_\star^{\mathrm{post}}
\mathbf H_{\mathrm{lik}}$ and
$\mathbf H_{\mathrm{prior}} \ll
(\alpha_\star^{\mathrm{post}})^2 \mathbf{JT}_{\mathrm{lik}}$, the $\alpha$-factors
inside the two log-determinants cancel and \eqref{eq:logZ-comb-def} reduces
to the clean additive form
\begin{equation}
\log Z_{\mathrm{comb}}
\;\approx\;
\log Z_{\mathrm{DCC}}
\;+\;
(\alpha_\star^{\mathrm{post}}-1)\,\log L(\hat\theta)
\;+\;
O\bigl(p\,\log\alpha_\star^{\mathrm{post}}\cdot
\|\mathbf H_{\mathrm{prior}}\|/\|\alpha_\star^{\mathrm{post}}
\mathbf H_{\mathrm{lik}}\|\bigr).
\label{eq:logZ-comb-leading}
\end{equation}
The leading-order net adjustment of $\log Z$ relative to the volume-corrected
$\log Z_{\mathrm{DCC}}$ is the single scalar $(\alpha_\star^{\mathrm{post}}-1)
\log L(\hat\theta)$, derived in Section~4.5.

\paragraph{Two orders of magnitude, one composition.}
The combined evidence couples corrections at two different orders in $T$:
the peak correction scales as $O(T)$ (since $\log L(\hat\theta) = O(T)$ under
stationarity and $(\alpha_\star^{\mathrm{post}}-1) = O(1)$), while the volume
correction inside $\log Z_{\mathrm{DCC}}$ scales as $O(1)$ (Section~9.1
intensive correction; Section~11.3 peak/volume resolution). The peak
correction therefore dominates the volume correction by a factor of $T$
in any persistent-misspecification regime. The intuition---``the peak
correction is the primary misspecification adjustment, the volume
correction is the secondary'' (introduced in Section~3)---is then a
quantitative ordering, not just a logical one.

\paragraph{When $\log Z_{\mathrm{comb}}$ matters for model comparison.}
For Bayes-factor comparisons between two models $A$ and $B$, the
combined-evidence Bayes-factor bias is
\begin{equation}
\Delta \log\mathrm{BF}_{AB}^{\mathrm{comb}}
\;=\;
\bigl(\alpha_\star^{A,\mathrm{post}}-1\bigr)\log L_A(\hat\theta_A)
\;-\;
\bigl(\alpha_\star^{B,\mathrm{post}}-1\bigr)\log L_B(\hat\theta_B)
\;+\;
\tfrac{1}{2}\log\det(\mathbf C_{\mathrm{post}}^{A})
\;-\;
\tfrac{1}{2}\log\det(\mathbf C_{\mathrm{post}}^{B}).
\label{eq:logBF-comb}
\end{equation}
The first two terms are the peak contributions; the last two are the volume
contributions of Sections~5--6. Section~9.2 discusses the practical
recommendation: report both contributions separately so reviewers can see
which one is driving model-comparison differences.

\paragraph{Evidence-interval extension.}
The evidence-interval framework of Section~10 below augments the triple
$(\log Z_{\mathrm{F}}, \log Z_{\mathrm{J}}, \log Z_{\mathrm{DCC}})$ to the
quadruple $(\log Z_{\mathrm{F}}, \log Z_{\mathrm{J}}, \log Z_{\mathrm{DCC}},
\log Z_{\mathrm{comb}})$. Section~10.4 discusses the pairwise differences in
the augmented framework.


\section*{8. Sign, Calibration, and Prior Cancellation}

\subsection*{8.1 Sign and magnitude}

The sign and magnitude of the correction are read off the scalar $\det(\mathbf C_{\mathrm{post}})$,
which is a multiplicative summary of the eigenvalues of $\mathbf C_{\mathrm{post}}$:
\begin{itemize}
\item $\det(\mathbf C_{\mathrm{post}}) > 1$ means net (multiplicative) variance inflation
across all eigen-directions of $\mathbf C_{\mathrm{post}}$, and implies $\Delta\log Z > 0$:
the corrected log-evidence is higher than the naive log-evidence, reflecting a larger
effective posterior volume.
\item $\det(\mathbf C_{\mathrm{post}}) < 1$ means net deflation across eigen-directions,
and implies $\Delta\log Z < 0$.
\item $\det(\mathbf C_{\mathrm{post}}) = 1$ implies $\Delta\log Z = 0$: naive and
DCC-corrected Laplace evidence coincide. This is a \emph{necessary but not sufficient}
condition for $\mathbf C_{\mathrm{post}} = \mathbf I$: a diagonal example such as
$\mathbf C_{\mathrm{post}} = \mathrm{diag}(2, 1/2)$ has unit determinant and zero
$\Delta\log Z$ yet is plainly miscalibrated. The scalar $\Delta\log Z$ therefore detects
miscalibration only up to per-direction cancellation; trace-based and per-eigenvalue
diagnostics from the main supplement (Sections~5.1 and~9) complement it.
\end{itemize}
For $\mathbf C_{\mathrm{post}} = c\,\mathbf I$ with $c > 0$, $\Delta\log Z = (p/2)\log c$,
scaling with dimension $p$ but not with sample size $T$.

\subsection*{8.2 Calibration under correct likelihood specification (population level)}

The defining property of the posterior correction is its behaviour at the population
level under correct likelihood specification, evaluated at the true parameter $\theta_0$.
By Section~6.1 of the main supplement,
\begin{equation}
\mathbf{JT}_{\mathrm{lik}}(\theta_0) = \mathbf H_{\mathrm{lik}}(\theta_0)
\;\;\Longrightarrow\;\;
\mathbf{JT}_{\mathrm{post}}(\theta_0) = \mathbf H_{\mathrm{post}}(\theta_0)
\;\;\Longrightarrow\;\;
\mathbf C_{\mathrm{post}}(\theta_0) = \mathbf I,
\label{eq:calibration-chain}
\end{equation}
the third implication holding for \emph{any} proper Gaussian prior. Substituting into
\eqref{eq:main-result},
\begin{equation}
\Delta\log Z
\;=\;
\tfrac{1}{2}\log\det(\mathbf I)
\;=\;
0
\qquad\text{at the population level under correct likelihood specification.}
\label{eq:DeltaZ-correct}
\end{equation}
The DCC and naive Bayesian Laplace evidences therefore agree at the population
level whenever the information identity holds, regardless of how informative the prior is.

\paragraph{Finite-sample caveat.}
The identity $\mathbf{JT}_{\mathrm{lik}} = \mathbf H_{\mathrm{lik}}$ is a population
identity at $\theta_0$ (main supp Sec~4.1); empirical plug-ins at $\hat\theta$ satisfy
it only up to $O_p(T^{-1/2})$ sampling fluctuation. The plug-in
$\mathbf C_{\mathrm{post}}(\hat\theta)$ therefore differs from $\mathbf I$ by $O_p(T^{-1/2})$
under correct spec, and the finite-sample shift
$\Delta\log Z = O_p(T^{-1})$ vanishes only asymptotically. The exact zero is a
population-level statement, not a finite-sample one.

\subsection*{8.3 Prior cancellation (population level under correct spec)}

The formula $\Delta\log Z = \tfrac{1}{2}\log\det(\mathbf C_{\mathrm{post}})$ depends on
$\mathbf H_{\mathrm{prior}}$ only implicitly through $\mathbf H_{\mathrm{post}}$ and
$\mathbf{JT}_{\mathrm{post}}$. At the population level under correct specification these
two matrices change in lockstep as the prior is rescaled (main supp Sec~3.1, final
paragraph), so $\mathbf C_{\mathrm{post}}$ remains identically $\mathbf I$ and the
correction \emph{cancels exactly} along the prior-strength axis. Concretely, define the
residual on the likelihood side
\begin{equation}
\Delta\mathbf{H}^{\mathrm{lik}}_{\mathrm{resid}}
\;:=\;
\mathbf{JT}_{\mathrm{lik}} - \mathbf H_{\mathrm{lik}}.
\label{eq:lik-residual}
\end{equation}
Then $\mathbf{JT}_{\mathrm{post}} - \mathbf H_{\mathrm{post}} = \Delta\mathbf{H}^{\mathrm{lik}}_{\mathrm{resid}}$
exactly (the prior contribution $\mathbf H_{\mathrm{prior}}$ cancels in the unwhitened
residual; main supp Sec~4.2), and
\begin{equation}
\mathbf C_{\mathrm{post}} - \mathbf I
\;=\;
\mathbf H_{\mathrm{post}}^{-1/2}\,
\Delta\mathbf{H}^{\mathrm{lik}}_{\mathrm{resid}}\,
\mathbf H_{\mathrm{post}}^{-1/2}.
\label{eq:Cpost-residual}
\end{equation}
The numerator $\Delta\mathbf{H}^{\mathrm{lik}}_{\mathrm{resid}}$ depends only on
likelihood quantities; the whitening factors, however, do depend on the prior, so the
determinant $\det(\mathbf C_{\mathrm{post}}) = \det(\mathbf{JT}_{\mathrm{post}})/\det(\mathbf H_{\mathrm{post}})$
is itself a prior-dependent readout under misspecification. The clean prior cancellation
in \eqref{eq:DeltaZ-correct} is a \emph{correct-spec phenomenon}, not a generic
prior-independence of the determinant; the strong-prior boundary in \S5.4 makes this
explicit.

\paragraph{Prior cancellation principle (population level under correct spec).}
\emph{Evaluated at $\theta_0$ under the population information identity, $\Delta\log Z = 0$
regardless of Gaussian-prior covariance. In this regime the posterior evidence correction
is a pure likelihood-misspecification diagnostic. Outside this regime---finite samples,
or misspecified models---the determinant readout is biased by the prior, as quantified
below.}

\subsection*{8.4 Strong-prior boundary: zero correction does not imply correct spec}

A second mechanism produces $\Delta\log Z \to 0$, independent of the population
information identity: a strong prior. By main supp Sec~4.2 (Eq.~\texttt{eq:posterior-c-residual-bound}),
\[
\|\mathbf C_{\mathrm{post}} - \mathbf I\|
\;\le\;
\lambda_{\min}(\mathbf H_{\mathrm{post}})^{-1}\,
\|\Delta\mathbf{H}^{\mathrm{lik}}_{\mathrm{resid}}\|,
\]
so as $\lambda_{\min}(\mathbf H_{\mathrm{prior}}) \to \infty$ the right-hand side vanishes
at rate $O(\lambda_{\min}(\mathbf H_{\mathrm{prior}})^{-1})$ \emph{even when the likelihood
is misspecified}, driving $\det(\mathbf C_{\mathrm{post}}) \to 1$ and hence $\Delta\log Z \to 0$.
Operationally, a near-zero $\Delta\log Z$ may indicate either (i) correct likelihood
specification at any prior, or (ii) a sufficiently strong prior that washes out the
likelihood-side residual. The scalar $\Delta\log Z$ alone cannot distinguish these two
causes; resolving the ambiguity requires either varying the prior strength (the strong-prior
limit cancels the residual uniformly across all data realisations) or inspecting per-eigenvalue
diagnostics of $\mathbf C_{\mathrm{lik}}$ on a retained subspace (main supp Sec~6.5 records
the same loss of absolute sensitivity).

\subsection*{8.5 Contrast with the likelihood evidence correction}

The likelihood evidence correction of \texttt{supplement\_evidence\_correction.tex} is
$\tfrac{1}{2}\log\det(\mathbf C)$ with $\mathbf C = \mathbf H_{\mathrm{lik}}^{-1/2}
\mathbf{JT}_{\mathrm{lik}} \mathbf H_{\mathrm{lik}}^{-1/2}$. When $\mathbf H_{\mathrm{lik}}$
is singular or near-singular (Section~8 of the Likelihood supplement) the likelihood-side
correction must be defined on a retained subspace, and finite-sample fluctuations along
weakly identified directions can leave $\mathbf C \neq \mathbf I$ even under correct
specification.

The posterior correction is structurally cleaner with respect to identifiability. Because
a proper Gaussian prior makes $\mathbf H_{\mathrm{post}} \succ 0$ on the full parameter
space (Section~2.1 of the main supplement, label \texttt{eq:dominance-condition}),
$\mathbf C_{\mathrm{post}}$ is well-defined as an SPD matrix without any subspace
projection. Under correct specification at $\theta_0$ it further reduces to $\mathbf I$
for any proper Gaussian prior (Section~6.1 of the main supplement). Identifiability is
thus handled by the prior; \emph{calibration} (whether $\mathbf C_{\mathrm{post}} = \mathbf I$)
is a separate question, answered by the diagnostic itself but---per \S5.4 above---only
in regimes where the prior is not so strong that the residual is suppressed.

\section*{9. Scaling with Sample Size and Bayes Factors}

\subsection*{9.1 Intensive correction}

Under stationarity and ergodicity of the per-interval scores with summable cross-interval
covariances (cf.\ main supp Secs~6.5, 7.3), $\mathbf H_{\mathrm{lik}}$ and
$\mathbf{JT}_{\mathrm{lik}}$ each scale linearly with $T$, the number of measurement intervals,
\begin{equation}
\mathbf H_{\mathrm{lik},T} \;\approx\; T\,\mathbf h_1,
\qquad
\mathbf{JT}_{\mathrm{lik},T} \;\approx\; T\,\boldsymbol\omega_1,
\qquad
\boldsymbol\omega_1 \;=\; \mathrm{Var}(s_1) + 2\sum_{k=1}^\infty \mathrm{Cov}(s_1, s_{1+k}),
\label{eq:lik-scaling}
\end{equation}
where $\mathbf h_1$ is the per-interval expected Fisher and $\boldsymbol\omega_1$ the
per-interval long-run score covariance; summability of the cross-interval covariances is
what makes $\boldsymbol\omega_1$ finite. The prior precision $\mathbf H_{\mathrm{prior}}$
is $O(1)$ in $T$. Hence
\begin{equation}
\mathbf H_{\mathrm{post},T}
\;=\;
T\,\mathbf h_1 + \mathbf H_{\mathrm{prior}}
\;=\;
T\bigl(\mathbf h_1 + T^{-1}\mathbf H_{\mathrm{prior}}\bigr),
\qquad
\mathbf{JT}_{\mathrm{post},T}
\;=\;
T\bigl(\boldsymbol\omega_1 + T^{-1}\mathbf H_{\mathrm{prior}}\bigr).
\label{eq:post-scaling}
\end{equation}
In the leading-order large-$T$ limit, the prior contributes a $T^{-1}$-suppressed
correction to both matrices on data-dominated directions, and
\begin{equation}
\mathbf C_{\mathrm{post},T}
\;\xrightarrow[T\to\infty]{}\;
\mathbf h_1^{-1/2}\,\boldsymbol\omega_1\,\mathbf h_1^{-1/2}
\;=:\;
\mathbf C_\infty,
\label{eq:Cpost-asymptote}
\end{equation}
an intensive limit independent of $T$ and of $\mathbf H_{\mathrm{prior}}$, provided
$\mathbf h_1 \succ 0$ (full data identifiability). Consequently
\begin{equation}
\Delta\log Z
\;=\;
\tfrac{1}{2}\log\det(\mathbf C_{\mathrm{post},T})
\;\xrightarrow[T\to\infty]{}\;
\tfrac{1}{2}\log\det(\mathbf C_\infty)
\;=\;
O(1),
\label{eq:DeltaZ-O1}
\end{equation}
an $O(1)$ shift that neither vanishes nor grows with $T$.

\subsection*{9.2 Bayes factors and decisive regime}

For two competing models $M_A$, $M_B$, the Bayes factor satisfies
\begin{equation}
\log\mathrm{BF}_{AB}
\;=\;
\log Z_A \;-\; \log Z_B,
\label{eq:BF-def}
\end{equation}
and the shift in the uncorrected (naive Bayesian Laplace) Bayes factor due to the DCC
correction is
\begin{equation}
\Delta\log\mathrm{BF}_{AB}
\;=\;
\tfrac{1}{2}\log\det(\mathbf C_{\mathrm{post},A})
\;-\;
\tfrac{1}{2}\log\det(\mathbf C_{\mathrm{post},B}),
\label{eq:BF-bias}
\end{equation}
an $O(1)$ quantity independent of $T$. By the sign of \eqref{eq:BF-bias}, the model whose
posterior distortion has the larger determinant gains relative weight under the DCC
correction; equivalently, the naive Bayes factor \emph{spuriously disfavours} the
more-distorted model by an amount that does not diminish with more data. The direction
of the bias is the opposite of the intuitive ``larger variance penalises a model'' reading,
because variance inflation enlarges the effective posterior volume entering $\log Z$.

\paragraph{Decisive regime.}
For most $T$, the $\log L(\hat\theta_A) - \log L(\hat\theta_B) = O(T)$ component of
$\log\mathrm{BF}_{AB}$ dominates and a fixed $O(1)$ correction is negligible. The
correction becomes decisive precisely in the regime where (i)~the $O(T)$ peak terms nearly
cancel (close models, or approximations of nearly equal fit), and (ii)~the prior is not
so strong that the likelihood-side residual is suppressed by the boundary mechanism of
\S5.4. This is the same operational regime identified for the likelihood-side correction
in \texttt{supplement\_evidence\_correction.tex}~\S5.3, restricted by the additional
prior-strength caveat.

\subsection*{9.3 Connection to posterior effective sample size}

Main supp Sec~9.5 defines a posterior effective sample size $T_{\mathrm{eff}}^{\mathrm{post}} = T/\kappa^{\mathrm{post}}$
via a spectral summary of $\mathbf R_{\mathrm{post}}$ (the distortion between $\mathbf{JT}_{\mathrm{post}}$
and $\mathbf{JS}_{\mathrm{post}}$). The evidence correction \eqref{eq:main-result}, in
contrast, is a scalar reduction of $\mathbf C_{\mathrm{post}}$ (the distortion between
$\mathbf{JT}_{\mathrm{post}}$ and $\mathbf H_{\mathrm{post}}$); the two statistics differ
in both their normaliser and their specific spectral summary. Using the multiplicative
split of Section~10 of the main supplement,
\begin{equation}
\det(\mathbf C_{\mathrm{post}})
\;=\;
\det(\mathbf C_{\mathrm{sample},\mathrm{post}})\,\det(\mathbf R_{\mathrm{post}}),
\label{eq:det-split}
\end{equation}
so the evidence correction decomposes as
\begin{equation}
\Delta\log Z
\;=\;
\tfrac{1}{2}\log\det(\mathbf C_{\mathrm{sample},\mathrm{post}})
\;+\;
\tfrac{1}{2}\log\det(\mathbf R_{\mathrm{post}}).
\label{eq:DeltaZ-split}
\end{equation}
The $\mathbf R_{\mathrm{post}}$ piece is the same matrix that drives
$T_{\mathrm{eff}}^{\mathrm{post}}$, but enters \eqref{eq:DeltaZ-split} as a log-determinant
rather than as the spectral summary $\kappa^{\mathrm{post}}$. When $\det(\mathbf C_{\mathrm{sample},\mathrm{post}}) = 1$
(within-interval calibration in the determinant sense), the evidence correction and the
$T_{\mathrm{eff}}^{\mathrm{post}}$ statistic track the same matrix $\mathbf R_{\mathrm{post}}$
via different scalar reductions; in general they remain distinct statistics.

\section*{10. The Evidence Interval Framework}

A practical advantage of the posterior framework is that three concurrent Laplace
estimators of $\log Z$ can be reported jointly, their relative spread being exactly
$|\log\det(\mathbf C_{\mathrm{post}})|$. The first is a derived Laplace approximation,
the third is the DCC correction derived above, and the second is a heuristic comparand
introduced explicitly to symmetrise the interval and to expose the calibration uncertainty
as a scalar.

\subsection*{10.1 Three Laplace-style estimators}

We define
\begin{align}
\log Z_F
\;&:=\;
\log L(\hat\theta) + \log\pi(\hat\theta)
+ \tfrac{p}{2}\log(2\pi)
- \tfrac{1}{2}\log\det(\mathbf H_{\mathrm{post}}),
\label{eq:ZF-def}\\[4pt]
\log Z_J
\;&:=\;
\log L(\hat\theta) + \log\pi(\hat\theta)
+ \tfrac{p}{2}\log(2\pi)
- \tfrac{1}{2}\log\det(\mathbf{JT}_{\mathrm{post}}),
\label{eq:ZJ-def}\\[4pt]
\log Z_{\mathrm{DCC}}
\;&:=\;
\log L(\hat\theta) + \log\pi(\hat\theta)
+ \tfrac{p}{2}\log(2\pi)
- \log\det(\mathbf H_{\mathrm{post}})
+ \tfrac{1}{2}\log\det(\mathbf{JT}_{\mathrm{post}}).
\label{eq:ZDCC-def}
\end{align}
$\log Z_F$ is the naive Bayesian Laplace evidence using the curvature
$\mathbf H_{\mathrm{post}}$ ($\log Z_F \equiv \log Z_{\mathrm{naive}}$ of \eqref{eq:naive-Z}).
$\log Z_{\mathrm{DCC}}$ is the DCC-corrected estimator obtained from \eqref{eq:corrected-Z}
via \eqref{eq:logdet-expansion}; both are derived from Laplace quadratic expansions of
$\log Z$ at $\hat\theta$.

$\log Z_J$, by contrast, is \emph{not} a derived Laplace approximation: the quadratic
expansion of \S2 identifies the local Gaussian precision with $\mathbf H_{\mathrm{post}}$,
not with $\mathbf{JT}_{\mathrm{post}}$. It is the formal substitution
$\mathbf H_{\mathrm{post}} \to \mathbf{JT}_{\mathrm{post}}$ in \eqref{eq:naive-Z}, included
solely to make the algebraic interval below symmetric about $\log Z_F$ and to expose the
calibration uncertainty $|\log\det(\mathbf C_{\mathrm{post}})|$ as a single scalar. We
call $\log Z_J$ the \emph{total-fluctuation comparand} (a J-only Laplace-style estimator);
the width of the interval $[\log Z_J, \log Z_{\mathrm{DCC}}]$ is then the algebraic gap
between this heuristic and the derived DCC estimate, not an error-theoretic confidence
band.

\subsection*{10.2 Pairwise differences}

A direct subtraction using \eqref{eq:detCpost-identity} gives
\begin{align}
\log Z_{\mathrm{DCC}} - \log Z_F
\;&=\;
-\tfrac{1}{2}\log\det(\mathbf H_{\mathrm{post}})
+ \tfrac{1}{2}\log\det(\mathbf{JT}_{\mathrm{post}})
\;=\;
+\tfrac{1}{2}\log\det(\mathbf C_{\mathrm{post}}),
\label{eq:ZDCC-minus-ZF}\\[4pt]
\log Z_J \;-\; \log Z_F
\;&=\;
-\tfrac{1}{2}\log\det(\mathbf{JT}_{\mathrm{post}})
+ \tfrac{1}{2}\log\det(\mathbf H_{\mathrm{post}})
\;=\;
-\tfrac{1}{2}\log\det(\mathbf C_{\mathrm{post}}),
\label{eq:ZJ-minus-ZF}\\[4pt]
\log Z_{\mathrm{DCC}} - \log Z_J
\;&=\;
-\log\det(\mathbf H_{\mathrm{post}})
+ \log\det(\mathbf{JT}_{\mathrm{post}})
\;=\;
+\log\det(\mathbf C_{\mathrm{post}}).
\label{eq:ZDCC-minus-ZJ}
\end{align}
All three pairwise differences are linear in the single scalar
$\log\det(\mathbf C_{\mathrm{post}})$:
\begin{equation}
\log Z_F,\;\log Z_J,\;\log Z_{\mathrm{DCC}}
\;\text{collapse to one point iff}\;
\log\det(\mathbf C_{\mathrm{post}}) = 0,
\label{eq:collapse-iff}
\end{equation}
and spread by $|\log\det(\mathbf C_{\mathrm{post}})|$ otherwise (the extreme pair is
$\log Z_J$ vs $\log Z_{\mathrm{DCC}}$).

\subsection*{10.3 Reporting the interval}

We propose to report the triple
\begin{equation}
\bigl(\log Z_F,\;\log Z_J,\;\log Z_{\mathrm{DCC}}\bigr)
\label{eq:triple}
\end{equation}
together with the interval
\begin{equation}
\bigl[\,
\min\{\log Z_F,\log Z_J,\log Z_{\mathrm{DCC}}\},
\;
\max\{\log Z_F,\log Z_J,\log Z_{\mathrm{DCC}}\}
\,\bigr]
\;=\;
\begin{cases}
[\log Z_J,\;\log Z_{\mathrm{DCC}}] & \text{if } \log\det(\mathbf C_{\mathrm{post}}) > 0,\\[4pt]
[\log Z_{\mathrm{DCC}},\;\log Z_J] & \text{if } \log\det(\mathbf C_{\mathrm{post}}) < 0,\\[4pt]
\{\log Z_F\} & \text{if } \log\det(\mathbf C_{\mathrm{post}}) = 0,
\end{cases}
\label{eq:Z-interval}
\end{equation}
with $\log Z_F$ always lying at the midpoint between $\log Z_J$ and $\log Z_{\mathrm{DCC}}$
(from \eqref{eq:ZDCC-minus-ZF}--\eqref{eq:ZJ-minus-ZF}). The width
\begin{equation}
W_Z
\;:=\;
|\log\det(\mathbf C_{\mathrm{post}})|
\label{eq:W-def}
\end{equation}
is a scalar calibration uncertainty: $W_Z = 0$ iff $\det(\mathbf C_{\mathrm{post}}) = 1$
(a necessary but not sufficient condition for $\mathbf C_{\mathrm{post}} = \mathbf I$;
per-eigenvalue cancellation can produce $W_Z = 0$ in genuinely miscalibrated cases), and
grows monotonically in net log-determinant deviation from unity. The interval framework
therefore provides a single report of $\log Z$ that exposes its determinant-level
calibration uncertainty due to information distortion, complementing the usual Monte Carlo
uncertainty of the Laplace approximation itself and the per-eigenvalue diagnostics of
Sections~5 and~9 of the main supplement.

\subsection*{10.4 Bayes factor under the interval framework}

For model comparison, applying \eqref{eq:BF-def} entry-wise gives three Bayes factor
estimators $\log\mathrm{BF}_{AB,F}$, $\log\mathrm{BF}_{AB,J}$, $\log\mathrm{BF}_{AB,\mathrm{DCC}}$.
From \eqref{eq:ZDCC-minus-ZJ},
\begin{equation}
\log\mathrm{BF}_{AB,\mathrm{DCC}}
\;-\;
\log\mathrm{BF}_{AB,J}
\;=\;
\log\det(\mathbf C_{\mathrm{post},A})
\;-\;
\log\det(\mathbf C_{\mathrm{post},B})
\;=\;
2\,(\Delta\log Z_A - \Delta\log Z_B),
\label{eq:BF-interval}
\end{equation}
so the spread between the DCC and J-only Bayes factors is twice the absolute difference
of the per-model $\Delta\log Z$ corrections. Note that this spread can vanish even when
both models are individually distorted (when $\det(\mathbf C_{\mathrm{post},A}) = \det(\mathbf C_{\mathrm{post},B})$)
or substantially exceed twice the per-model $|\Delta\log Z|$ when the two corrections have
opposite signs. A decisive Bayes factor that exceeds this spread in absolute value is
robust to determinant-level calibration uncertainty; one inside the spread is not, and the
user should either explicitly model the misspecification or accept the Bayes factor as
undetermined.

\subsection*{10.5 Position of $\log Z_{\mathrm{comb}}$ relative to the interval}

The combined evidence $\log Z_{\mathrm{comb}}$ of Section~7 is not a fourth member
of the interval $[\log Z_J,\log Z_{\mathrm{DCC}}]$: the interval reflects only the
volume-term ambiguity ($\mathbf H_{\mathrm{post}}$ vs.\ $\mathbf{JT}_{\mathrm{post}}$
in the determinant), while $\log Z_{\mathrm{comb}}$ additionally re-weights the peak
term by the trace-matched learning rate
$\alpha_\star^{\mathrm{post}} = p/\mathrm{tr}(\mathbf C_{\mathrm{post}})$
(Section~4). Decomposing \eqref{eq:logZ-comb-clean} in the obvious way,
\begin{equation}
\log Z_{\mathrm{comb}}
\;=\;
\underbrace{\log Z_{\mathrm{DCC}}}_{\text{volume-corrected, }O(1)\text{ in }T}
\;+\;
\underbrace{\bigl(\alpha_\star^{\mathrm{post}}-1\bigr)\log L(\hat\theta)}_{\text{peak correction, }O(T)\text{ under persistent misspec}},
\label{eq:logZcomb-vs-interval}
\end{equation}
so $\log Z_{\mathrm{comb}}$ shifts the whole interval rigidly by the peak term while
preserving its width $W_Z = |\log\det(\mathbf C_{\mathrm{post}})|$. The reported
quadruple $(\log Z_F,\log Z_J,\log Z_{\mathrm{DCC}},\log Z_{\mathrm{comb}})$ then
contains two operationally distinct readouts: the interval width measures volume
calibration uncertainty (a function only of $\det(\mathbf C_{\mathrm{post}})$),
while the gap
$\log Z_{\mathrm{comb}} - \log Z_{\mathrm{DCC}} = (\alpha_\star^{\mathrm{post}}-1)\log L(\hat\theta)$
measures the peak-term reweighting (a function of $\mathrm{tr}(\mathbf C_{\mathrm{post}})$).

For Bayes factors, the combined-evidence version $\log B_{AB}^{\mathrm{comb}}$ inherits
the volume-correction Bayes-factor spread of \eqref{eq:BF-interval} unchanged and adds
the per-model peak-correction difference
$(\alpha_{\star,A}^{\mathrm{post}}-1)\log L_A(\hat\theta_A) -
 (\alpha_{\star,B}^{\mathrm{post}}-1)\log L_B(\hat\theta_B)$.
Under persistent misspecification this peak-correction difference is $O(T)$ and
typically dominates the $O(1)$ volume-correction spread; under
correct specification both contributions are calibrated and the combined Bayes
factor reduces to the naive Bayes factor up to $O_p(T^{-1/2})$ noise.

\section*{11. The Apparent Paradox: Logarithmic Variance Inflation}

The compactness of \eqref{eq:main-result} can be misleading. As in
\texttt{supplement\_evidence\_correction.tex}~\S6, we make explicit why a correction that
``should'' be huge---a fold-doubling of posterior variance corresponds to a halved effective
sample size---turns out to be merely an $O(1)$ shift in $\log Z$.

\subsection*{11.1 Variance inflation as effective sample loss}

The PID acts as a variance-inflation factor for the MAP estimator (Section~9 of the main
supplement); CI widths scale as $\sqrt{[\mathbf C_{\mathrm{post}}]_{ii}}$, so
$\mathbf C_{\mathrm{post}} = 2\mathbf I$ doubles posterior variance and widens intervals
by $\sqrt{2}$. By the usual variance--information duality, this is equivalent to a halved
\emph{evidence-equivalent effective sample size}, which we denote
\begin{equation}
N_{\mathrm{eff}}^{Z}
\;:=\;
T\,\det(\mathbf C_{\mathrm{post}})^{-1/p}
\label{eq:Neff-from-Cpost}
\end{equation}
in the scalar-inflation case. We emphasise that $N_{\mathrm{eff}}^{Z}$ is a
$\mathbf C_{\mathrm{post}}$-based scalar reduction relevant to the evidence shift
$\Delta\log Z$; it is \emph{distinct} from the posterior effective sample size
$T_{\mathrm{eff}}^{\mathrm{post}} = T/\kappa^{\mathrm{post}}$ defined in main supp
Sec~9.5 via a spectral summary of $\mathbf R_{\mathrm{post}}$, and the two coincide only
when $\det(\mathbf C_{\mathrm{sample},\mathrm{post}}) = 1$ and the determinant of
$\mathbf R_{\mathrm{post}}$ matches its $\kappa^{\mathrm{post}}$ summary (cf.\
\eqref{eq:det-split}). Halving $N_{\mathrm{eff}}^{Z}$ feels like a major loss of information.

\subsection*{11.2 Yet the log-evidence shift is small}

Substituting $\mathbf C_{\mathrm{post}} = 2\mathbf I$ for $p = 6$ into \eqref{eq:main-result},
\begin{equation}
\Delta\log Z
\;=\;
\tfrac{p}{2}\log 2
\;\approx\;
2.08\ \mathrm{nats}
\;\approx\;
3.0\ \mathrm{bits},
\label{eq:numerical-shift}
\end{equation}
corresponding to a Bayes factor of $\approx 8$. Notable for model comparison but
negligible compared to typical $O(T)$ log-evidence values.

\subsection*{11.3 Resolution: peak vs.\ volume}

Decompose the Laplace evidence into a peak and a volume contribution,
\begin{equation}
\log Z
\;\approx\;
\underbrace{\log L(\hat\theta)}_{\text{peak} \;=\; O(T)}
\;+\;
\underbrace{\log\pi(\hat\theta)}_{\text{prior peak} \;=\; O(1)}
\;+\;
\tfrac{p}{2}\log(2\pi)
\;+\;
\underbrace{\tfrac{1}{2}\log\det(\Sigma)}_{\text{volume}}.
\label{eq:peak-volume-split}
\end{equation}
The peak term $\log L(\hat\theta)$ depends on how well the model fits at its best
parameters and grows linearly in $T$. The prior peak $\log\pi(\hat\theta)$ is fixed in $T$
(once $\hat\theta$ stabilises) since the prior is a fixed observation. The volume term
further splits.

The naive volume contribution is $-\tfrac{1}{2}\log\det(\mathbf H_{\mathrm{post}})$. Let
$p_{\mathrm{eff}} \le p$ be the number of data-dominated directions of $\mathbf H_{\mathrm{post}}$
(those whose $\mathbf h_1$-eigenvalues are bounded away from zero); the remaining
$p - p_{\mathrm{eff}}$ directions are prior-dominated with $\mathbf H_{\mathrm{post}}$-eigenvalues
$O(1)$. Then
\begin{equation}
-\tfrac{1}{2}\log\det(\mathbf H_{\mathrm{post}})
\;=\;
-\tfrac{p_{\mathrm{eff}}}{2}\log T \;+\; O(1)
\qquad \text{(BIC-like penalty on data-dominated directions)}.
\label{eq:BIC-penalty}
\end{equation}
When $p_{\mathrm{eff}} = p$ this reduces to the textbook BIC penalty; when some directions
are prior-dominated (main supp Sec~2.1 dominance condition; main supp Sec~6.5 retained
eigenspace), the penalty contracts to $p_{\mathrm{eff}}$. The DCC correction adds
$\tfrac{1}{2}\log\det(\mathbf C_{\mathrm{post}})$, which by \eqref{eq:Cpost-asymptote} is
$O(1)$ on data-dominated directions:
\begin{equation}
\Delta\log Z
\;=\;
\tfrac{1}{2}\log\det(\mathbf C_{\mathrm{post}})
\;=\;
O(1)
\qquad \text{as } T\to\infty.
\label{eq:correction-O1}
\end{equation}
Variance inflation is therefore a \emph{second-order} effect on $\log Z$: it modifies the
shape of the posterior around the peak but not the peak itself, and the volume term
contains both the BIC-like $O(\log T)$ penalty (from $\mathbf H_{\mathrm{post}}$) and the
$O(1)$ distortion correction (from $\mathbf C_{\mathrm{post}}$). Doubling effective sample
size halves $\det(\mathbf C_{\mathrm{post}})^{1/p}$ but adds only $\tfrac{p}{2}\log 2$ to
$\log Z$.

\subsection*{11.4 Multiplicative in probability, additive in log-probability}

Evidence accumulates multiplicatively over independent samples,
$Z_T \approx \prod_{t=1}^T p(y_t \mid \mathrm{history}_t)$, and so $\log Z_T$ is additive
in $T$. Doubling the data approximately doubles the $O(T)$ peak contribution to $\log Z_T$;
the volume term contributes only $-(p_{\mathrm{eff}}/2)\log T + O(1)$ and is subdominant.
Doubling the posterior variance multiplies $Z_T$ by $2^{p/2}$ and adds the constant
$(p/2)\log 2$ to $\log Z_T$. These are fundamentally different scalings:
\begin{itemize}
\item \emph{More data}: additive in $\log Z$, $O(T)$ leading contribution (peak term).
\item \emph{Variance inflation}: additive in $\log Z$, $O(1)$ contribution (PID).
\end{itemize}
For large $T$ the data term dominates, and the DCC correction is a fine adjustment; for
model comparison where the $O(T)$ peak terms approximately cancel, the $O(1)$ correction
becomes the leading discriminator. As in the likelihood case
(\texttt{supplement\_evidence\_correction.tex}~\S6.4), the same correction object is
``small'' or ``decisive'' depending entirely on what it is being compared against.

\section*{12. Summary}

\begin{itemize}
\item \textbf{Two corrections, one diagnostic.} Likelihood misspecification distorts
\emph{both} the peak and the volume term of the naive Laplace evidence; each requires
its own correction, and both extract complementary scalar summaries of the same
$\mathrm{eig}(\mathbf C_{\mathrm{post}})$ spectrum. The peak correction
(Section~4) re-weights $\log L(\hat\theta)$ via
$\alpha_\star^{\mathrm{post}} = p/\mathrm{tr}(\mathbf C_{\mathrm{post}})$ (arithmetic
mean); the volume correction (Sections~5--6) shifts the volume term by
$\tfrac{1}{2}\log\det(\mathbf C_{\mathrm{post}})$ (geometric mean). The combined
evidence $\log Z_{\mathrm{comb}}$ of Section~7 incorporates both.
\item \textbf{Volume-correction central result.} Replacing the naive Bayesian Laplace
covariance $\mathbf H_{\mathrm{post}}^{-1}$ with the sandwich $\Sigma_{\mathrm{DCC},\mathrm{post}}$
shifts the log-evidence by
\[
\Delta\log Z
\;=\;
\tfrac{1}{2}\log\det(\mathbf C_{\mathrm{post}}),
\]
the direct posterior analogue of the likelihood evidence correction
(\texttt{supplement\_evidence\_correction.tex}).
\item \textbf{Peak-correction central result.} The combined evidence decomposes as
$\log Z_{\mathrm{comb}} = \log Z_{\mathrm{DCC}} +
(\alpha_\star^{\mathrm{post}}-1)\log L(\hat\theta)$. The peak term is $O(T)$ under
persistent misspecification and dominates the $O(1)$ volume term; under correct
specification $\alpha_\star^{\mathrm{post}} \to 1$ and the peak correction vanishes.
\item \textbf{Prior cancellation (population level under correct spec).}
At $\theta_0$ under the population information identity, $\mathbf C_{\mathrm{post}} = \mathbf I$
for \emph{any} proper Gaussian prior covariance, so $\Delta\log Z = 0$ regardless of prior
strength. In this regime the posterior correction is a pure likelihood-misspecification
diagnostic. Finite-sample plug-ins satisfy this only up to $O_p(T^{-1/2})$ fluctuation,
and the determinant readout is biased toward unity in the strong-prior limit even under
misspecification ($\|\mathbf C_{\mathrm{post}}-\mathbf I\| \to 0$ at rate
$O(\lambda_{\min}(\mathbf H_{\mathrm{prior}})^{-1})$). A zero $\Delta\log Z$ therefore
does not by itself imply correct specification.
\item \textbf{Intensive scaling.} $\Delta\log Z$ is $O(1)$ in $T$ under stationarity,
ergodicity, and summability of cross-interval score covariances. It neither grows nor
vanishes with sample size and is decisive only in the regime where (i)~$O(T)$ peak terms
nearly cancel and (ii)~the prior is sufficiently weak that the likelihood-side residual
is not suppressed.
\item \textbf{Evidence interval framework.} The three Laplace-style estimators
$(\log Z_F,\,\log Z_J,\,\log Z_{\mathrm{DCC}})$---two derived, one ($\log Z_J$) included
as a J-only comparand---collapse to a single point when $\log\det(\mathbf C_{\mathrm{post}}) = 0$
and spread by $|\log\det(\mathbf C_{\mathrm{post}})|$ otherwise. The width
$W_Z = |\log\det(\mathbf C_{\mathrm{post}})|$ is a determinant-level scalar calibration
uncertainty (necessary but not sufficient for $\mathbf C_{\mathrm{post}} = \mathbf I$);
per-eigenvalue diagnostics complement it.
\item \textbf{Subspace projection avoided.} A proper Gaussian prior makes
$\mathbf H_{\mathrm{post}} \succ 0$ on the full parameter space (Section~2.1 of the main
supplement, label \texttt{eq:dominance-condition}), so $\mathbf C_{\mathrm{post}}$ is
well-defined as an SPD matrix and $\Delta\log Z$ is real-valued without the
retained-eigenspace machinery required for the singular-Fisher likelihood case
(\texttt{supplement\_evidence\_correction.tex}~\S5). The dominance condition itself remains
a precondition for finite-difference Hessians with indefinite spectrum.
\item \textbf{Effective sample sizes are distinct.} The
$\mathbf C_{\mathrm{post}}$-based $N_{\mathrm{eff}}^{Z} := T\det(\mathbf C_{\mathrm{post}})^{-1/p}$
governs the evidence shift; the $\mathbf R_{\mathrm{post}}$-based
$T_{\mathrm{eff}}^{\mathrm{post}} = T/\kappa^{\mathrm{post}}$ of main supp Sec~9.5 governs
score-vs-Fisher concentration. They coincide only when within-interval calibration holds
in the determinant sense.
\item \textbf{Apparent paradox resolved.} Variance inflation, which feels like a major
loss of effective sample size, contributes only an $O(1)$ additive shift to $\log Z$
because it lives entirely in the volume term of the Laplace split, not in the $O(T)$ peak
term. Bayesian model comparison senses this $O(1)$ shift only where the $O(T)$ contributions
nearly cancel.
\end{itemize}

\end{document}
